\documentclass[11pt]{article}

\usepackage[margin=1in]{geometry}
\usepackage[T1]{fontenc}
\usepackage[utf8]{inputenc}
\usepackage{lmodern}
\usepackage{microtype}
\usepackage{amsmath,amssymb,amsthm}
\usepackage{mathtools}
\usepackage{hyperref}
\usepackage{booktabs}
\usepackage{enumitem}
\usepackage{listings}
\usepackage{xcolor}

\hypersetup{
  colorlinks=true,
  linkcolor=blue!50!black,
  citecolor=blue!50!black,
  urlcolor=blue!50!black
}

\lstdefinelanguage{Lean4}{
  morekeywords={
    import,namespace,end,open,section,noncomputable,abbrev,
    def,theorem,lemma,where,by,fun,let,in,have,show,exact,
    intro,apply,simpa,calc
  },
  sensitive=true,
  morecomment=[l]{--},
  morecomment=[s]{/-}{-/},
  morestring=[b]"
}

\lstset{
  language=Lean4,
  basicstyle=\ttfamily\small,
  keywordstyle=\bfseries\color{blue!50!black},
  commentstyle=\itshape\color{green!40!black},
  stringstyle=\color{red!50!black},
  showstringspaces=false,
  breaklines=true,
  frame=single,
  framerule=0.3pt,
  xleftmargin=0.5em,
  xrightmargin=0.5em,
  keepspaces=true,
  columns=fullflexible,
  literate=
    {¬}{{$\neg$}}1
    {∃}{{$\exists$}}1
    {→}{{$\to$}}1
    {≤}{{$\leq$}}1
}

\newtheorem{theorem}{Theorem}[section]
\newtheorem{lemma}[theorem]{Lemma}
\theoremstyle{remark}
\newtheorem{remark}[theorem]{Remark}

\title{P vs NP Proof Attempts:\\
Crux-First and Steelman Lean Notes on a Weakness-Based Route}
\author{}
\date{\today}

\begin{document}
\maketitle

\begin{abstract}
This note records the current Lean work around a claimed proof of
$P \neq NP$ via ``quantale weakness and geometric complexity.''  The project now
proceeds on two tracks in parallel.

The \emph{blocker} track asks where the argument must fail unless an additional theorem
is supplied.  The \emph{optimistic} track assumes the high-level strategy may work and
builds the missing background theory from the ground up so that any genuine repair can
plug into a machine-checked scaffold.

The current Lean work proves concrete obstructions.  First, a decoder that is merely
``local on a bounded-degree neighborhood of radius $\Theta(\log m)$'' is \emph{not} thereby
polylogarithmic in size: the visible neighborhood already has polynomial size in $m$.
Second, if one does not prove an additional compression theorem, then the class of all local
rules on that neighborhood is doubly exponential in the neighborhood size, so it cannot be
uniformly represented by a polylogarithmic code budget.

Before the switching story even starts, the sampler already carries another burden.  The
manuscript fixes the VV width at \(k=c_1\log m\), but the cited \(\Omega(1/m)\) uniqueness
probability comes from averaging over many scales of \(k\).  The new Lean obstruction packages the
fixed-width variance bound: if the retained-solution count has mean \(\mu>1\) and variance at most
\(\mu\), then exact isolation occurs on at most a \(\mu/(\mu-1)^2\) fraction of hash choices.
So fixed logarithmic width can only isolate often when the pre-hash solution set is already
essentially polynomial in size.

And even the manuscript's repeated \emph{pairwise-independent columns} language cannot shoulder the
hashing argument by itself.  A separate Lean countermodel shows that pairwise-independent columns do
not imply universal hashing: one can have pairwise-uniform columns together with a deterministic
three-column linear dependency, so a nonzero difference hashes to zero with probability \(1\).

Nor can the manuscript's \(\delta\)-biased right-hand side rescue the hash moments by itself.
Another Lean module shows that once a candidate's label map is genuinely uniform over the hash-label
space, its average hit mass against \emph{any} finite probability weight on labels is exactly the
uniform value \(1/|L|\); and if a pair of candidates has jointly uniform labels, then the weighted
average joint-hit mass is exactly \(1/|L|^2\).  So any useful effect of the \(\delta\)-biased
right-hand side must come from failure of the underlying label maps to be genuinely uniform, not
from the bias itself.

And even exact pre-conditioning uniformity is not enough by itself.  A further Lean countermodel
shows that conditioning on a label-dependent event can collapse a perfectly uniform label map to a
point mass.  So any appeal to pre-conditioning \(2\)-universality or uniform labels must be
supplemented by a separate theorem about what survives the manuscript's later conditioning on
uniqueness.

Third, even one natural repair to the calibration step has a sharp failure mode: if an
involution flips the target bit while preserving the retained local features, then every
classifier using only those invariant features has exactly \(1/2\) average accuracy.  So one
cannot repair the sign-flip calibration bug merely by dropping the involution-sensitive part of
the local input and hoping marginalization will preserve a domination claim.

Even before that charitable repair, the manuscript's exact local input formula already has a hard
consistency problem: it defines the post-switch input as \(u=(z,a_i,b)\) while the involution
\(T_i\) toggles \(b\) by \(a_i\).  So full-input invariance can hold only on the zero-column case
\(a_i=0\).

The new bridge module sharpens that fork: on the exact same \(T_i\)-action, the reduced projection
\((z,a_i)\) is invariant and therefore carries zero symmetry-respecting signed signal, while any
advantage obtained by also retaining \(b\) is bounded by the mass on nonzero VV columns, since
those are exactly the points where \(b\) resolves an involution pair.

Sharper still, the same involution hypotheses force every retained feature fiber to contain
exactly half target-\(1\) and half target-\(0\) points.  So even soft scores or conditional-mean
estimators on those invariant features have zero margin on every feature value.

And this is not just a uniform-counting artifact: any involution-invariant finite weighting still
assigns equal total mass to the target-\(1\) and target-\(0\) parts of every retained feature
fiber.  So non-uniform but symmetry-respecting sampling does not restore a usable margin.

In fact, any soft score computed only from those involution-invariant retained features has zero
weighted signed correlation with the target.  So even a richer invariant scoring statistic carries
no usable directional signal before genuine symmetry breaking is introduced.

More sharply, even when the weighting is not exactly involution-invariant, the entire invariant-score
signal is still carried only by the antisymmetric part of that weighting.  So this rescue route now
needs a real theorem that substantial useful mass lives in the weight asymmetry itself.

Likewise, if one augments the invariant local view by a side channel, then any advantage gained by
classifying on the pair \((u,v)\) is still bounded by the total mass where the side channel actually
separates an orbit point from its involution partner.

Equivalently, any claimed domination gap over chance now forces a matching amount of orbit-resolving
mass.  So one cannot even state a serious advantage theorem without simultaneously proving a serious
resolution theorem.

More generally, keeping a controlled amount of non-invariant information only helps on the slice
where it really breaks the symmetry.  Any residual slice on which the retained features are still
involution-invariant remains exactly at chance under symmetry-respecting weights.

Quantitatively, any excess over chance is bounded by the weight of the
symmetry-broken region.  So a proof that keeps only a little non-invariant information must also
show that almost all relevant mass actually lies outside the unresolved symmetric slice.

Fourth, the natural ``short global program'' rescue also fails on its own: if the switched input
really has only \(n = O(\log m)\) visible bits, then an arbitrary local rule on those bits can be
hard-wired by a truth table of size \(2^n = \mathrm{poly}(m)\).  So polynomial-size decoders do
not by themselves force any small subclass inside the full local-rule space.

Fifth, the most creative metagraph-style rescue faces its own exactness barrier: if a quotient
state must capture the full interface behavior of a \(k\)-bit local gadget, then there are
\(2^{k2^k}\) such boundary transducers.  At binary-log width this already forces at least linear
code size in \(m\), so a small exact metagraph/message quotient cannot appear by default.

Sixth, even the tree-message route does not follow from tree-likeness alone: a fixed perfect
binary tree architecture of depth \(n\) already realizes every Boolean rule on \(n\) visible bits,
just by changing its leaf labels.

At the same time, the optimistic Lean work now isolates the precise positive interface one
would need for the switching lemma to go through: a \emph{small explicitly encoded
hypothesis family} whose realized classifier class is bounded by its code budget, together
with a perspective-relative notion of accessible input surface inspired by the existing
Hyperseed formalization.

These results do not by themselves disprove the whole paper.  They do isolate an essential gap:
the switching/locality step must be strengthened by a real theorem showing that the specific
switched predictors lie in an exceptionally small subclass, not merely in the class of all local
functions.
\end{abstract}

\tableofcontents

\section{Context}

The proof strategy under review is the October 2025 manuscript
\emph{``P != NP: A Non-Relativizing Proof via Quantale Weakness and Geometric Complexity.''}
At a high level, the paper proposes the following chain:
\begin{enumerate}[leftmargin=1.5em]
  \item Construct a masked random Unique-SAT ensemble.
  \item Show that every short polynomial-time decoder can be wrapped into a predictor that is
  local on a constant fraction of blocks.
  \item Use neutrality and sparsification to show local predictors succeed only with exponentially
  small probability across many independent blocks.
  \item Convert that small-success statement into a lower bound on time-bounded Kolmogorov
  complexity.
  \item Contradict the $P=NP$ upper bound obtained from Unique-SAT self-reduction.
\end{enumerate}

The current crux-first effort does not attempt the full chain at once.  Instead it attacks the
most fragile inference in the middle:

\begin{quote}
``Switched predictor is local at logarithmic radius'' $\Longrightarrow$
``switched predictor belongs to a polylogarithmic hypothesis class.''
\end{quote}

That implication is not automatic.  It needs either a structural theorem about the switched
predictor or a direct circuit/description-length bound.  The Lean work below shows why.

\section{Current Lean Files}

The current Lean modules under
\texttt{Mettapedia/Computability/PNP/}:
\begin{center}
\begin{tabular}{ll}
\toprule
\textbf{File} & \textbf{Role} \\
\midrule
\texttt{HypothesisClass.lean} & optimistic code-indexed hypothesis families and Hyperseed-visible inputs \\
\texttt{LocalityObstruction.lean} & Log-radius locality is larger than polylog size \\
\texttt{CompressionObstruction.lean} & Arbitrary local rules do not admit short uniform encodings \\
\texttt{FixedWidthIsolationObstruction.lean} & fixed-width hashing isolates only on a small fraction once the expected survivor count exceeds \(1\) \\
\texttt{PairwiseSurvivorMoments.lean} & pairwise first and factorial-second moment bounds imply the fixed-width variance budget \\
\texttt{PairwiseCandidateBridge.lean} & candidate-wise \(0/1\) hits plus pairwise off-diagonal bounds instantiate the survivor-moment model \\
\texttt{PairwiseColumnsObstruction.lean} & pairwise-independent columns still need not form a universal hash family \\
\texttt{RhsBiasIrrelevance.lean} & once label maps are uniform, any finite RHS bias averages to the same single-hit and pair-hit moments \\
\texttt{TwoUniversalRhsIrrelevance.lean} & exact finite \(2\)-universality of the hash family makes RHS bias irrelevant to the first two hit moments \\
\texttt{ConditioningObstruction.lean} & conditioning can collapse a perfectly uniform label map to a point mass \\
\texttt{SymmetrizationObstruction.lean} & unbiased majority does not beat Bayes without a margin \\
\texttt{PostSwitchInputObstruction.lean} & exact \(u=(z,a_i,b)\) is \(T_i\)-invariant iff \(a_i=0\) \\
\texttt{PostSwitchForkObstruction.lean} & exact \(T_i\)-fork: invariant \((z,a_i)\) has zero signal, \(b\)-side-channel gain needs nonzero-column mass \\
\texttt{VisiblePostSwitchSurface.lean} & isolates the exact visible surface, invariant projection, and fork-visible projection \\
\texttt{ExactSwitchedFamily.lean} & isolates the exact switched-family compression target and ERM reuse theorem \\
\texttt{ExactVisibleCompressionObstruction.lean} & if the exact family spans all exact visible rules, no sub-surface-cardinality code budget can cover it \\
\texttt{OrbitNeutralityObstruction.lean} & involution-invariant feature collapse forces exact half accuracy \\
\texttt{FiberNeutralityObstruction.lean} & invariant feature fibers have exact zero conditional margin \\
\texttt{WeightedFiberNeutralityObstruction.lean} & symmetry-respecting weights still give zero conditional margin \\
\texttt{InvariantScoreObstruction.lean} & invariant soft scores have zero weighted signed signal \\
\texttt{WeightAsymmetryObstruction.lean} & invariant-score signal comes only from antisymmetric weight mass \\
\texttt{SideChannelResolutionObstruction.lean} & side channels help only on orbit-resolving mass \\
\texttt{ResolutionDemandObstruction.lean} & any domination gap forces matching orbit-resolving mass \\
\texttt{ResidualSymmetryObstruction.lean} & any unresolved symmetric slice still stays at chance \\
\texttt{AsymmetryBudgetObstruction.lean} & advantage is bounded by the symmetry-broken mass \\
\texttt{ShortProgramObstruction.lean} & polynomial-size programs already realize all log-width local rules \\
\texttt{MetagraphObstruction.lean} & exact metagraph/interface quotients have state explosion \\
\texttt{TreeMessageObstruction.lean} & fixed tree architectures already realize all local rules \\
\bottomrule
\end{tabular}
\end{center}

These compile in the current \texttt{mettapedia} workspace via \texttt{lake env lean}
(with selected files also built to \texttt{.olean} artifacts).

\section{Two Parallel Programs}

\subsection{The Blocker Program}

The blocker program is adversarial.  It asks:
\begin{enumerate}[leftmargin=1.5em]
  \item Is the visible radius already too large for any naive \texttt{poly(log m)} claim?
  \item Is the full class of local rules too large for a short description?
  \item Can the manuscript's fixed \(k=c_1\log m\) VV layer really achieve \(\Omega(1/m)\)
  uniqueness probability on large solution sets, or does fixed-width isolation already collapse
  once the expected survivor count is bigger than \(1\)?
  \item Is the manuscript's weaker \emph{pairwise-independent columns} language alone enough to
  justify universal hashing, or can higher-order linear dependencies still make a nonzero
  difference hash to zero deterministically?
  \item If the hash labels are already uniform or pairwise-uniform, can the \(\delta\)-biased
  right-hand side still change the first and pair hit moments at all?
  \item Even if the pre-conditioning label map is uniform, can later conditioning on uniqueness
  still collapse the local label distribution completely?
  \item Does majority symmetrization fail whenever the local conditional mean has zero margin?
  \item Is the exact post-switch input \(u=(z,a_i,b)\) really preserved by the paper's involution,
  or only in the zero-column case \(a_i=0\)?
  \item Once that exact fork is made explicit, does the charitable repair reduce either to zero
  invariant-feature signal on \((z,a_i)\) or to advantage bounded by the mass on nonzero VV columns?
  \item Does the obvious ``drop the bad coordinate'' calibration patch collapse any feature-only
  classifier back to exact neutrality?
  \item Does the softer ``estimate the conditional mean on invariant features'' repair still have
  exact zero margin on every retained feature fiber?
  \item Does non-uniform but involution-symmetric weighting avoid that zero-margin collapse?
  \item Can an arbitrary invariant soft score recover directional signal, or does it also have zero
  signed correlation with the target?
  \item If the weighting is not exactly symmetric, does invariant-score signal come only from the
  antisymmetric part of that weighting?
  \item If one keeps a non-invariant side channel, is total advantage still bounded by the mass
  where that side channel actually resolves involution pairs?
  \item Conversely, does any claimed domination advantage force a comparably large orbit-resolving
  mass for that side channel?
  \item If one keeps some controlled non-invariant information, do residual unresolved slices still
  remain exact-chance unless their mass is proved negligible?
  \item Can a small symmetry-breaking slice still yield a large global advantage, or is the excess
  over chance quantitatively bounded by the symmetry-broken mass?
  \item Does ``short global decoder'' alone still allow the full local-rule class at log-width?
  \item Does an exact finite-state metagraph/message quotient on a log-width interface already
  require too many states?
  \item Does tree-likeness alone force a small exact message/interpreter class?
\end{enumerate}

The current Lean answer to all twenty-one is ``yes.''  Therefore the published proof cannot be
accepted \emph{as written} unless it supplies additional structural theorems.

\subsection{The Optimistic Program}

The optimistic program asks what would be enough to save the route.
The current working answer is:
\begin{quote}
locality alone is not enough, but \emph{locality plus explicit compressibility} may be.
\end{quote}

Concretely, the switching theorem would need to produce not merely a local rule on a
radius-$O(\log m)$ neighborhood, but a family of predictors realized by a code space of
size at most \(\mathrm{polylog}(m)\) bits.  In that case the realized hypothesis class is
small, and standard finite-class learning bounds become plausible again.

The new module \texttt{HypothesisClass.lean} formalizes exactly this interface:
\begin{enumerate}[leftmargin=1.5em]
  \item an \emph{encoded family} is a decoder from a finite code space into classifiers;
  \item the realized class has cardinality at most the size of the code space;
  \item therefore an $s$-bit family realizes at most $2^s$ classifiers;
  \item no such family can surject onto the full local-rule class once $s < 2^n$.
\end{enumerate}

This is the correct positive abstraction: if this switching argument is salvageable, it must
land in one of these small encoded subclasses rather than in the full local-rule space.

\subsection{Hyperseed Perspective Surface}

The Hyperseed formalization in \texttt{Mettapedia/Hyperseed/} already provides a useful
observer-relative interface:
\begin{itemize}[leftmargin=1.5em]
  \item a \emph{perspective} determines what is observable,
  \item a budget determines what is affordable,
  \item a guard determines what is admitted.
\end{itemize}

The optimistic PNP background layer reuses this by treating a decoder's \emph{available
input surface} as a subtype of the Hyperseed \texttt{availableRegion}.  This does not solve
switching.  It does provide a clean ground-up semantics for what a bounded observer can even
use as input, without baking in a bespoke notion of locality for every later argument.

\section{Optimistic Background Theorem}

The basic positive theorem is elementary but indispensable:

\begin{theorem}[Code budget bounds realized hypothesis class]
Let \(H\) be a family of Boolean classifiers realized by decoding an \(s\)-bit code:
\[
  \mathrm{decode} : \{0,1\}^s \to \mathcal{F}.
\]
Then the realized class satisfies
\[
  |\mathrm{range}(\mathrm{decode})| \le 2^s.
\]
\end{theorem}

\noindent
This is formalized in \texttt{HypothesisClass.lean} as
\texttt{BitEncodedClassifierFamily.card\_realized\_le}.

The contrapositive is equally important:

\begin{theorem}[Short codes cannot cover the full local-rule class]
If \(s < 2^n\), then no \(s\)-bit decoder can surject onto the full set of Boolean rules on
\(n\) visible bits.
\end{theorem}

\noindent
This is formalized as
\texttt{not\_surjective\_decode\_to\_fullLocalRule\_of\_lt}.

These statements are not the final theorem.  They are the right \emph{background theory} for
the optimistic route: they tell us exactly what shape of compression theorem the switching
step must deliver.

\section{Obstruction I: Log-Radius Locality Is Already Polynomial}

Suppose a factor graph has bounded degree $d \ge 2$, and a predictor is allowed to inspect a
radius-$r$ neighborhood around one target bit.  A rough but unavoidable upper/lower size model is
\[
  \text{visible nodes} \asymp d^r.
\]
If $r = \Theta(\log m)$, then this is
\[
  d^{c \log m} = m^{c \log d},
\]
which is polynomial in $m$, not polylogarithmic in $m$.

The Lean module \texttt{LocalityObstruction.lean} records this in several exact forms.

\begin{theorem}[Binary-log lower bound]
For every $d,m \in \mathbb{N}$ with $d \ge 2$,
\[
  m \le d^{\log_2 m + 1}.
\]
\end{theorem}

\noindent
This is formalized as the theorem
\texttt{inputSize\_le\_neighborhoodSize\_at\_binaryLogRadius}.

\begin{theorem}[Power-of-two specialization]
For every $d,n \in \mathbb{N}$ with $d \ge 2$,
\[
  2^n \le d^{\log_2(2^n)} = d^n.
\]
\end{theorem}

\noindent
This appears as
\texttt{powerOfTwo\_le\_neighborhoodSize\_at\_exactBinaryLogRadius}.

\begin{theorem}[Asymptotic dominance]
Fix $k \in \mathbb{N}$ and real constants $d>1$, $c>0$.  Then
\[
  (\log x)^k = o\!\left(d^{c \log x}\right)
\]
as $x \to \infty$.
\end{theorem}

\noindent
This is formalized as
\texttt{polylog\_isLittleO\_logRadiusNeighborhood}, and the corresponding negative
Big-O statement appears as
\texttt{logRadiusNeighborhood\_not\_isBigO\_polylog}.

\subsection{Interpretation}

This does \emph{not} say that every local rule requires polynomial-size circuits.
It says something weaker but still decisive: locality at logarithmic radius, by itself, does not
justify a \texttt{poly(log m)} complexity claim.  The radius already exposes a polynomial-sized
input window.

So if the paper wants a tiny post-switch hypothesis class, it must prove a \emph{compression
theorem} about the specific local rules that arise after switching.  That is a separate and
nontrivial burden.

\section{Obstruction Ib: Fixed-Width VV Isolation Does Not Follow from the Cited Lemma}

The manuscript's main ensemble fixes
\[
  k=c_1\log m
\]
in the VV layer.  But the cited isolation lemma proving uniqueness probability \(\Omega(1/m)\)
chooses \(k\) \emph{uniformly across all scales} \(0,\dots,m-1\).  Those are not the same claim.

The new Lean module \texttt{FixedWidthIsolationObstruction.lean} packages the exact finite bound
one needs at fixed width.

\subsection{Lean theorem}

\begin{theorem}[Variance-budget bound for exact isolation]
Let \(Y\) be the number of candidates retained by a random hash choice, on a finite hash space.
If
\[
  \mathbb E[Y]=\mu>1
  \qquad\text{and}\qquad
  \mathrm{Var}(Y)\le \mu,
\]
then
\[
  \Pr[Y=1]\ \le\ \frac{\mu}{(\mu-1)^2}.
\]
\end{theorem}

\noindent
This is formalized as
\texttt{exactOneProb\_le\_of\_variance\_budget}.

The new companion module \texttt{PairwiseSurvivorMoments.lean} formalizes the abstract bridge from
pairwise-survivor moments to this variance budget.  In that file, a finite model satisfying
\[
  \mathbb E[Y]=\mu
  \qquad\text{and}\qquad
  \mathbb E[Y(Y-1)]\le \mu^2
\]
is shown to obey \(\mathrm{Var}(Y)\le \mu\), and therefore the same exact-isolation upper bound.
This separates the remaining manuscript obligation cleanly: one still has to prove that the actual
fixed-width VV sampler instantiates those pairwise moments.

The newer bridge module \texttt{PairwiseCandidateBridge.lean} takes one more step toward the actual
VV picture.  It proves that if the retained-candidate indicators are \(0/1\)-valued, all have the
same marginal hit rate \(p\), and distinct candidates satisfy the pairwise upper bound
\[
  \mathbb E[X_a X_b] \le p^2,
  \qquad a\neq b,
\]
then the total survivor count satisfies the exact abstract interface of
\texttt{PairwiseSurvivorMomentModel}.  Concretely:
\begin{itemize}[leftmargin=1.5em]
  \item \texttt{candidateHitCount\_sum} proves the first moment formula;
  \item \texttt{candidateHitCount\_factorialSecondMoment\_le} proves the factorial second-moment
  bound;
  \item \texttt{pairwiseCandidateMomentModel} packages these into the survivor-moment model used by
  the fixed-width isolation theorem.
\end{itemize}

So the remaining gap has narrowed again.  To recover the manuscript's fixed-width uniqueness claim,
one now needs an end-to-end theorem that the actual VV hash family satisfies these candidate-wise
marginal and pairwise bounds at the fixed logarithmic width used later in the paper.

\subsection{Why this hits the manuscript}

In the source text, the supporting VV lemma states the classical bound
\[
  \Pr[|S\cap h^{-1}(0^k)|=1]\ge c/m
\]
only after choosing \(k\) uniformly from \(0,\dots,m-1\).  But the actual ensemble later fixes
\(k=c_1\log m\) and then still says isolation succeeds with \(\Omega(1/m)\) probability.

That gap is not cosmetic.  Under the standard pairwise-independent VV heuristic, if a formula has
\(N\) satisfying assignments and the fixed hash width is \(k\), then the retained-count random
variable has mean
\[
  \mu = N/2^k
\]
and variance at most \(\mu\).  The Lean theorem above therefore yields
\[
  \Pr[\text{exact isolation}]
  \ \lesssim\
  \frac{1}{\mu}
  \ =\
  \frac{2^k}{N},
\]
up to a universal constant once \(\mu\) is bounded away from \(1\).

So with fixed \(k=O(\log m)\), the sampler can isolate with noticeable probability only when the
pre-hash solution set is already \(N=O(2^k)=\mathrm{poly}(m)\).  The manuscript, however, only
cites a much weaker upper bound of the form \(N\le 2^{\alpha m}\), which still permits an
exponential solution set.  On such instances, fixed logarithmic width gives exponentially small
exact-isolation probability rather than \(\Omega(1/m)\).

\subsection{Consequence}

This means the paper now needs a real theorem of one of the following forms:
\begin{enumerate}[leftmargin=1.5em]
  \item on the relevant random-3CNF mass, the satisfiable formulas entering the VV stage already
  have only \(\mathrm{poly}(m)\) satisfying assignments; or
  \item the ensemble does not actually use fixed \(k=c_1\log m\), and the later local-input and
  finite-alphabet claims must be rewritten accordingly.
\end{enumerate}

Without one of those repairs, the promised block distribution \(D_m\) is not supported by the
cited VV argument.

\section{Obstruction Ic: Pairwise-Independent Columns Do Not Imply Universal Hashing}

The ensemble definition samples \(A\) from ``any \(2\)-universal distribution with
pairwise-independent columns,'' and the later robustness discussion repeatedly foregrounds the
\emph{pairwise-independent columns} part.  The new Lean module
\texttt{PairwiseColumnsObstruction.lean} shows that this weaker condition is not enough on its own.

\subsection{Lean theorem}

\begin{theorem}[Pairwise columns need not be universal]
There exists a one-bit linear hash family with three columns such that:
\begin{enumerate}[leftmargin=1.5em]
  \item each pair of columns is jointly uniform over \(\mathbb{F}_2^2\); but
  \item some nonzero input difference hashes to \(0\) on \emph{every} seed.
\end{enumerate}
Hence pairwise-independent columns do not by themselves imply the universal-hashing bound required
for VV isolation.
\end{theorem}

\noindent
This is formalized by the bijectivity theorems
\texttt{col12\_bijective}, \texttt{col13\_bijective}, and \texttt{col23\_bijective}, together with
\texttt{threeColumnHash\_badDifference} and
\texttt{threeColumnHash\_not\_universal}.

\subsection{Why this matters}

This does \emph{not} refute a hash family that is genuinely \(2\)-universal.  It does refute a
common informal shortcut: one cannot replace real universal hashing by mere pairwise independence
of the columns.  Higher-order linear dependencies can survive that weaker condition and completely
destroy the needed collision bound.

So the manuscript now owes a sharper statement here too:
\begin{enumerate}[leftmargin=1.5em]
  \item either prove that the actual fixed-width VV family is genuinely \(2\)-universal in the
  sense needed for the candidate-wise moment bridge; or
  \item stop appealing to pairwise-independent columns as if that weaker property alone carried the
  isolation argument.
\end{enumerate}

\section{Obstruction Id: RHS Bias Is Irrelevant Once The Label Maps Are Uniform}

The manuscript samples the VV right-hand side \(b\) from a \(\delta\)-biased source and suggests
that this robustness can be folded into the collision/isolation analysis.  The new Lean module
\texttt{RhsBiasIrrelevance.lean} isolates the exact point where that can and cannot matter.

\subsection{Lean theorems}

\begin{theorem}[Uniform labels erase single-hit bias]
Let \(L\) be a finite label space with any probability weight \(w\) on \(L\), and suppose a finite
seed space \(S\) splits as \(S \cong \{1,\dots,d\} \times L\) with \(d>0\), so that a candidate's
label map is exactly the projection onto \(L\).  Then
\[
  \frac{1}{|S|}\sum_{s \in S} w(f(s)) = \frac{1}{|L|}.
\]
\end{theorem}

\noindent
This is formalized by \texttt{average\_mass\_of\_uniform\_labels}.

\begin{theorem}[Jointly uniform labels erase pair-hit bias]
Under the same hypotheses, if a pair of candidates has jointly uniform labels
\(S \cong \{1,\dots,d\} \times L \times L\) with \(d>0\), then
\[
  \frac{1}{|S|}
  \sum_{s \in S} \mathbf{1}_{f(s)=g(s)}\, w(f(s))
  = \frac{1}{|L|^2}.
\]
\end{theorem}

\noindent
This is formalized by \texttt{average\_joint\_mass\_of\_uniform\_label\_pairs}.

\subsection{Why this matters}

This does not prove that the manuscript's actual fixed-width VV family has perfectly uniform label
maps.  It proves the sharper conditional statement: \emph{if} the candidate-wise and pairwise label
maps are genuinely uniform, then the choice of a \(\delta\)-biased right-hand-side distribution is
irrelevant to the first and pair-hit moments.  The averaged values are exactly the same as under a
uniform right-hand side.

So the manuscript cannot get extra collision-control power from \(\delta\)-bias alone.  Any real
effect of the right-hand side must come from a failure of the underlying label maps to be uniform in
the first place.  This narrows the burden again:
\begin{enumerate}[leftmargin=1.5em]
  \item either prove that the actual fixed-width family's label maps are \emph{not} uniform in a
  way that helps the needed moment bounds; or
  \item stop treating \(\delta\)-bias itself as if it repaired the missing universal-hashing step.
\end{enumerate}

The new bridge module \texttt{TwoUniversalRhsIrrelevance.lean} sharpens this one step further in
the manuscript's own language: if the actual \(A\)-family is genuinely finite \(2\)-universal, then
the weighted single-hit and pair-hit averages are exactly the uniform values for \emph{every}
finite right-hand-side distribution \(w\).  So on the first two VV moments, the entire burden sits
on the \(A\)-family, not on the \(\delta\)-biased choice of \(b\).

But even this sharper bridge is only a \emph{pre-conditioning} statement.  The new module
\texttt{ConditioningObstruction.lean} shows that this is the end of what one can get for free:
conditioning a perfectly uniform label map on a single label fiber makes the conditional label
distribution a point mass.  So no argument may pass from pre-conditioning uniformity to the
distribution conditioned on uniqueness without an additional theorem controlling that conditioning
step.

\section{Obstruction II: The Full Local Rule Class Is Enormous}

Let $n$ be the number of visible Boolean features after switching.  Then the unconstrained class of
local predictors is simply
\[
  \{0,1\}^{\{0,1\}^n},
\]
the set of all Boolean functions on $n$ bits.  Its size is
\[
  2^{2^n}.
\]
Therefore any \emph{uniform} code space that represents all such rules injectively must have at
least $2^n$ bits.

This is formalized in \texttt{CompressionObstruction.lean}.

\begin{theorem}[Counting bound for local rules]
Let \texttt{LocalRule n} denote the type of Boolean functions on \texttt{n} visible bits, and
\texttt{BitCode s} the type of \texttt{s}-bit codewords.  If
\[
  s < 2^n,
\]
then there is no injective map
\[
  \texttt{LocalRule n} \to \texttt{BitCode s}.
\]
\end{theorem}

\noindent
This is the theorem \texttt{no\_injective\_bitCode\_of\_lt}.

\begin{theorem}[Binary-log neighborhood corollary]
If $d \ge 2$ and
\[
  s < 2^m,
\]
then there is no injective map
\[
  \texttt{LocalRule}\!\left(d^{\log_2 m + 1}\right) \to \texttt{BitCode s}.
\]
\end{theorem}

\noindent
This is the theorem
\texttt{no\_uniform\_code\_below\_expInput\_for\_binaryLogNeighborhood}.

\begin{theorem}[Size lower bound on the full class]
For $d \ge 2$,
\[
  2^{2^m} \le
  \left|\texttt{LocalRule}\!\left(d^{\log_2 m + 1}\right)\right|.
\]
\end{theorem}

\noindent
This is the theorem
\texttt{expExpInput\_le\_card\_localRule\_binaryLogNeighborhood}.

\subsection{Interpretation}

Again, this does not say that the switched rule class \emph{is} the full local rule class.  It says
that without a theorem proving otherwise, one cannot pass from
``local on $n$ visible bits'' to ``describable by only $\mathrm{polylog}(m)$ bits.''

That gap is exactly where the current program remains vulnerable.

\section{What These Results Break}

The present Lean work targets the following inference schema:
\[
  \text{short decoder}
  \Longrightarrow
  \text{switched local predictor}
  \Longrightarrow
  \text{tiny ACC$^0$/ERM hypothesis}.
\]

The first arrow is the paper's switching claim.  The second arrow is where the present obstruction
lands.

\begin{remark}
If the switched predictor is only known to be some arbitrary function of a log-radius neighborhood,
then:
\begin{enumerate}[leftmargin=1.5em]
  \item the neighborhood already has polynomial size in the ambient parameter, and
  \item the class of all such local rules is far too large to admit a polylogarithmic uniform code.
\end{enumerate}
\end{remark}

Therefore a successful proof attempt must supply one of the following:
\begin{enumerate}[leftmargin=1.5em]
  \item a theorem that the switched predictors belong to a very special subclass
  (for example, a family with explicit circuit/description bound), or
  \item a direct compression algorithm showing that the relevant local rules can be encoded in
  \(\mathrm{polylog}(m)\) bits for reasons stronger than locality.
\end{enumerate}

Without such a theorem, the proof's route from locality to tiny hypothesis class is unsupported.

\section{Obstruction III: Zero-Margin Symmetrization Cannot Concentrate}

The symmetrization route in the paper makes a stronger claim than the two counting obstructions
above.  It says, roughly, that after taking a polylogarithmic number of promise-preserving sign
flips and then taking a majority vote, the resulting surrogate label agrees with the Bayes rule on
the local $\sigma$-field for all but a negligible fraction of blocks.

There is a direct obstruction here when the conditional surrogate mean lands at exactly
\(\tfrac12\).

\subsection{Lean counting lemma}

The Lean module \texttt{SymmetrizationObstruction.lean} formalizes the following fact.

\begin{theorem}[Odd-majority is exactly balanced on unbiased bits]
Let \(s\) be odd.  Among the \(2^s\) bit-vectors of length \(s\), exactly half have strict majority
\texttt{true} and half have strict majority \texttt{false}.
\end{theorem}

\noindent
The key formal statements are:
\begin{itemize}[leftmargin=1.5em]
  \item \texttt{majority\_compl}: complementing all bits flips the majority bit when \(s\) is odd;
  \item \texttt{card\_majority\_true\_eq\_card\_majority\_false}: the majority-\texttt{true} and
  majority-\texttt{false} classes have equal cardinality;
  \item \texttt{majority\_tie\_break\_not\_high\_probability}: if a Bayes rule breaks a
  \(p=\tfrac12\) tie toward \texttt{true}, then majority on \(s\) unbiased samples matches that rule
  on only half the sample space.
\end{itemize}

\subsection{Application to the Current Symmetrization Wrapper}

Now consider the paper's symmetrized surrogate label for a fixed block \(j\) and bit \(i\).  Take
the simplest possible decoder \(P\): the decoder that always outputs \(0\).

Under a sign flip \(\sigma\), the paper's own back-mapping rule gives
\[
  Y^{(r)}_{j,i} = \langle a_{j,i}, \sigma^{(r)} \rangle
\]
over \(\mathbb{F}_2\), because the decoder output is \(0\) and the back-map xors in the VV shift.

If \(a_{j,i} \neq 0\), then \(\langle a_{j,i}, \sigma \rangle\) is an unbiased bit under uniform random
\(\sigma\).  Hence the sequence
\[
  Y^{(1)}_{j,i}, \dots, Y^{(s)}_{j,i}
\]
is a fair sample of surrogate bits, and for odd \(s\) the majority vote is exactly balanced.

But the paper's Bayes rule is written with threshold \(1[p \ge \tfrac12]\), which at \(p=\tfrac12\)
chooses the value \(1\).  Therefore, in this zero-margin case, the majority vote matches the stated
Bayes rule on only half the seed choices, not with probability \(1 - m^{-\Omega(1)}\).

\subsection{Consequence}

This is materially stronger than the earlier ``local does not imply simple'' objection.  It says:

\begin{quote}
The concentration-to-Bayes step in the symmetrization argument is false unless an additional
margin theorem is proved.
\end{quote}

In other words, one must show that for the relevant blocks and bits,
\[
  \left| p_{j,i} - \tfrac12 \right|
\]
is bounded below by more than the sampling noise scale.  Without such a margin, majority voting
cannot recover the Bayes rule uniformly, even for the constant-zero decoder.

\section{Obstruction IIIb: The Exact Post-Switch Input Is Not \(T_i\)-Invariant}

Before moving to a charitable repair, the manuscript's own exact definitions already create a
direct problem.

Section 3.2 defines the promise-preserving involution
\[
  T_i(F^h, A, b) = (F^{\tau_i h}, A, b \oplus a_i),
\]
while Section 3.3 defines the post-switch local input
\[
  u_i(\Phi) = (z(F^h), a_i, b).
\]
Appendix A.5 then says this same involution preserves \(u=(z,a_i,b)\).  The new Lean module
\texttt{PostSwitchInputObstruction.lean} formalizes the exact contradiction point.

\subsection{Lean theorem}

\begin{theorem}[Full-input invariance holds only on zero VV columns]
Let \(u = (z,a,b)\) be the exact post-switch local input, and let \(T\) act on it by
\[
  T(z,a,b) = (z,a,b \oplus a).
\]
Then
\[
  T(u) = u
  \qquad\Longleftrightarrow\qquad
  a = 0.
\]
Equivalently, if the VV column \(a\) is nonzero, then the involution necessarily changes the
right-hand side coordinate \(b\).
\end{theorem}

\noindent
The key formal statements are:
\begin{itemize}[leftmargin=1.5em]
  \item \texttt{vvToggle\_eq\_self\_iff\_zero}: \(b \oplus a = b\) iff \(a=0\);
  \item \texttt{tiInputMap\_eq\_self\_iff\_zeroColumn}: the exact local input is \(T_i\)-invariant
  iff the VV column is zero;
  \item \texttt{invariantProjection\_tiInputMap}: the reduced projection \((z,a)\) is always
  preserved;
  \item \texttt{tiInputMap\_changes\_b\_of\_nonzeroColumn}: if \(a \neq 0\), then \(b\) really
  changes.
\end{itemize}

\subsection{Why this matters}

This is not a soft complaint about omitted detail.  It is a direct incompatibility between the
paper's own formulas:
\begin{enumerate}[leftmargin=1.5em]
  \item the exact local input keeps \(b\),
  \item the involution toggles \(b\) by \(a_i\),
  \item yet the calibration appendix says the involution preserves \(u=(z,a_i,b)\).
\end{enumerate}

That can only be true on the special case \(a_i=0\).  So the calibration argument cannot be right
as written on the full local input.

\subsection{Consequence}

This forces the proof into a fork:
\begin{enumerate}[leftmargin=1.5em]
  \item either keep the full input \((z,a_i,b)\), in which case the involution is not feature
  preserving and a new calibration theorem is needed;
  \item or drop the changing coordinate \(b\), in which case one lands in the invariant-feature
  regime attacked next.
\end{enumerate}

\section{Obstruction IIIc: The Exact Manuscript Fork Already Lands in the Symmetry Budget}

The previous section exposed the formula-level inconsistency in the manuscript as written.  The new
Lean module \texttt{PostSwitchForkObstruction.lean} shows that the most charitable exact repair is
already constrained by the later symmetry-budget files.

\subsection{Lean theorems}

\begin{theorem}[The exact \(T_i\)-action is an involution and \(b\) resolves exactly the nonzero columns]
On the exact local input \(u=(z,a,b)\) with
\[
  T(z,a,b) = (z,a,b \oplus a),
\]
the map \(T\) is an involution.  Moreover,
\[
  b(Tu) \neq b(u)
  \qquad\Longleftrightarrow\qquad
  a \neq 0.
\]
Equivalently, the \(b\)-coordinate distinguishes an involution pair exactly on the nonzero-column
slice.
\end{theorem}

\begin{theorem}[The exact fork]
Let the target bit flip under the exact involution \(T\), and let the weighting be \(T\)-invariant.
Then:
\begin{enumerate}[leftmargin=1.5em]
  \item every soft score depending only on the reduced invariant projection \((z,a)\) has zero
  weighted signed correlation with the target;
  \item every classifier using \(((z,a),b)\) has success advantage over chance bounded by the
  total weight of the nonzero-column slice.
\end{enumerate}
\end{theorem}

\noindent
The key formal statements are:
\begin{itemize}[leftmargin=1.5em]
  \item \texttt{tiInputMap\_involutive}: the manuscript's exact local update really is an
  involution;
  \item \texttt{resolvedBySideChannel\_tiInputMap\_b\_iff\_nonzeroColumn}: the \(b\)-side channel
  resolves an orbit pair exactly when \(a \neq 0\);
  \item \texttt{weighted\_signedScore\_sum\_eq\_zero\_on\_invariantProjection}: the reduced exact
  input \((z,a)\) has zero signed signal under symmetry-respecting weights;
  \item \texttt{doubledAdvantage\_invariantProjection\_with\_b\_le\_nonzeroColumnMass}: any
  classifier on \(((z,a),b)\) pays for its domination gap with nonzero-column mass.
\end{itemize}

\subsection{Why this matters}

This turns the manuscript fork into a direct quantitative obligation.

If one drops \(b\), then the surviving exact input is already inside the invariant-feature regime:
no symmetry-respecting signed signal remains.  If one keeps \(b\), then the proof is no longer
allowed to say merely that ``a small side channel helps.''  It must prove that the resulting
domination gap is supported by a sufficiently large nonzero-column slice in the actual weighted
domain.

So this exact bridge does not merely identify a local inconsistency.  It shows that the natural
repair lands immediately in the later neutrality and resolution-budget obstructions.

\section{Obstruction IV: Involution-Invariant Feature Collapse Forces Exact Neutrality}

Once the exact full-input inconsistency is noticed, the most natural charitable repair is to drop
the involution-sensitive coordinate \(b\) and keep only the involution-fixed part of the local view,
say a feature map of the form
\[
  u' = (z, a_i).
\]

This repair looks attractive because it restores feature invariance.  The new Lean module
\texttt{OrbitNeutralityObstruction.lean} shows why this move cannot, by itself, support a
domination theorem.

\subsection{Lean theorem}

\begin{theorem}[Feature-preserving involutions force exact half accuracy]
Let \(\alpha\) be a finite sample space, let \(\tau : \alpha \to \alpha\) be an involution, let
\(u : \alpha \to U\) be a feature map, and let \(y : \alpha \to \{0,1\}\) be the target bit.
Assume
\[
  u(\tau x) = u(x)
  \qquad\text{and}\qquad
  y(\tau x) = 1 - y(x)
\]
for every \(x \in \alpha\).  Then for every classifier \(h : U \to \{0,1\}\),
\[
  2 \cdot \#\{x : h(u(x)) = y(x)\} = |\alpha|.
\]
Equivalently, the average accuracy of \(h \circ u\) is exactly \(1/2\).
\end{theorem}

\noindent
The key formal statements are:
\begin{itemize}[leftmargin=1.5em]
  \item \texttt{correct\_under\_pair\_iff\_incorrect}: correctness at \(x\) is equivalent to
  incorrectness at \(\tau x\);
  \item \texttt{card\_correct\_eq\_card\_incorrect}: the correct and incorrect classes have equal
  cardinality;
  \item \texttt{two\_mul\_card\_correct\_eq\_card}: exactly half the finite sample space is
  classified correctly.
\end{itemize}

\subsection{Why this hits the charitable calibration patch}

Suppose the repair to the calibration bug is:
\begin{enumerate}[leftmargin=1.5em]
  \item remove the involution-sensitive coordinate \(b\) from the local input;
  \item prove neutrality by pairing each instance with its \(T_i\)-flipped partner;
  \item then learn or evaluate a classifier using only the reduced feature map \(u'\).
\end{enumerate}

If the paired instances really preserve \(u'\) and flip the target bit, then the theorem above
applies immediately.  Every feature-only classifier on \(u'\) has exact \(1/2\) average accuracy.
So this route cannot yield a strict domination statement, a high-probability Bayes recovery
statement, or any per-block success advantage above chance.

\subsection{Consequence}

This is a stronger objection than ``the paper needs more detail.''  It says:

\begin{quote}
If the calibration fix works by dropping the coordinate changed by the involution and then
classifying from the invariant remainder alone, that fix destroys any chance of beating
\(1/2\) accuracy on the paired finite domain.
\end{quote}

So a viable rescue must do something more subtle than mere feature collapse.  For example, it
would have to:
\begin{enumerate}[leftmargin=1.5em]
  \item keep some controlled non-invariant information,
  \item switch from hard classification to a softer statistic whose expectation matters,
  \item or show that the actual learning/evaluation domain is not the balanced finite orbit space
  modeled by the theorem above.
\end{enumerate}

\section{Obstruction IVb: Invariant Feature Fibers Have Zero Margin}

The most natural reply to Obstruction IV is:
\begin{quote}
Fine, hard classification on the invariant features fails.  But one can still estimate a
\emph{soft score} or conditional mean on those features and only use high-margin feature values.
\end{quote}

This is the right stronger target.  The new Lean module
\texttt{FiberNeutralityObstruction.lean} blocks that repair under the same involution
hypotheses.

\subsection{Lean theorem}

\begin{theorem}[Every invariant feature fiber is exactly balanced]
Let \(\alpha\) be a finite sample space, let \(\tau : \alpha \to \alpha\) be an involution, let
\(u : \alpha \to U\) be a feature map, and let \(y : \alpha \to \{0,1\}\) be the target bit.
Assume
\[
  u(\tau x) = u(x)
  \qquad\text{and}\qquad
  y(\tau x) = 1 - y(x)
\]
for every \(x \in \alpha\).  Then for every retained feature value \(v \in U\),
\[
  2 \cdot \#\{x : u(x)=v \text{ and } y(x)=1\}
  =
  \#\{x : u(x)=v\}.
\]
Equivalently, conditioned on any retained feature value \(v\), the target bit has exact mean
\(1/2\).
\end{theorem}

\noindent
The key formal statements are:
\begin{itemize}[leftmargin=1.5em]
  \item \texttt{featureFiberTrueEquivFalse}: on each fiber \(u^{-1}(v)\), the involution pairs
  the \texttt{true} points with the \texttt{false} points;
  \item \texttt{card\_featureFiberTrue\_eq\_card\_featureFiberFalse}: the two parts of each fiber
  have equal cardinality;
  \item \texttt{two\_mul\_card\_featureFiberTrue\_eq\_card\_featureFiber}: the \texttt{true}
  points occupy exactly half of each fiber.
\end{itemize}

\subsection{Why this hits the soft-score repair}

This is stronger than the hard-classifier obstruction.  It says the retained feature map does not
merely fail to support a good classifier; it fails to support \emph{any nonzero conditional
margin}.  On every retained feature value \(v\), the exact local average is already
\[
  \mathbb{E}[\,y \mid u=v\,] = \tfrac12.
\]

So if the repair route is:
\begin{enumerate}[leftmargin=1.5em]
  \item drop the involution-sensitive coordinate,
  \item keep only the invariant summary \(u'\),
  \item estimate a soft score or conditional mean from \(u'\),
  \item then use only ``high-margin'' \(u'\)-values,
\end{enumerate}
it fails immediately on the balanced finite orbit model.  There are no high-margin retained
feature values at all.

\subsection{Consequence}

This narrows the remaining rescue space again.  A successful patch now has to do more than
``replace hard labels by soft scores.''  It must either:
\begin{enumerate}[leftmargin=1.5em]
  \item retain some controlled non-invariant information,
  \item show that the actual averaging domain is not fiberwise balanced under the involution,
  \item or use a structure theorem stronger than feature-invariant orbit pairing.
\end{enumerate}

\section{Obstruction IVc: Symmetry-Respecting Weights Still Have Zero Margin}

Another natural reply is:
\begin{quote}
The real evaluation or training average is not uniform on the finite orbit space.  Once one uses
the correct non-uniform weighting, the invariant features may regain a useful soft margin.
\end{quote}

The new Lean module \texttt{WeightedFiberNeutralityObstruction.lean} blocks that repair whenever
the weighting itself respects the same involution symmetry.

\subsection{Lean theorem}

\begin{theorem}[Orbit-symmetric weights preserve zero fiber margin]
Let \(\alpha\) be a finite sample space, let \(\tau : \alpha \to \alpha\) be an involution, let
\(u : \alpha \to U\) be a retained feature map, let \(y : \alpha \to \{0,1\}\) be the target bit,
and let \(w : \alpha \to M\) be a weight into any additive commutative monoid.  Assume
\[
  u(\tau x) = u(x), \qquad y(\tau x) = 1-y(x), \qquad w(\tau x) = w(x)
\]
for every \(x \in \alpha\).  Then for every retained feature value \(v\),
\[
  \sum_{x : u(x)=v,\ y(x)=1} w(x)
  =
  \sum_{x : u(x)=v,\ y(x)=0} w(x).
\]
\end{theorem}

\noindent
The key formal statements are:
\begin{itemize}[leftmargin=1.5em]
  \item \texttt{weightedFeatureFiberTrueMass}: the weighted mass of the \texttt{true} part of one
  feature fiber;
  \item \texttt{weightedFeatureFiberFalseMass}: the weighted mass of the \texttt{false} part;
  \item \texttt{weightedFeatureFiberTrueMass\_eq\_weightedFeatureFiberFalseMass}: the exact balance
  theorem under involution-invariant weights.
\end{itemize}

\subsection{Why this closes the non-uniform averaging reply}

This matters because it removes a very common escape hatch.  One can no longer say:
\begin{quote}
The balanced finite orbit space was just a counting toy model.  The real proof uses a more
careful empirical or distributional weighting, so the invariant local statistic may still have a
nonzero expectation gap.
\end{quote}

Not if that weighting respects the same sign-flip involution.  The theorem shows that every such
weighting still gives exact zero conditional margin on every retained feature fiber.  In
particular, normalizing an involution-invariant finite weight into a probability distribution
still leaves
\[
  \mathbb{E}[\,y \mid u=v\,] = \tfrac12
\]
for every retained feature value \(v\) with nonzero mass.

\subsection{Consequence}

So the remaining escape routes are narrower again.  To revive a margin on the retained local
input, one must now do something stronger than:
\begin{enumerate}[leftmargin=1.5em]
  \item pass from hard labels to soft scores,
  \item pass from uniform counts to non-uniform but symmetry-respecting weights,
  \item or restrict attention to invariant-feature fibers alone.
\end{enumerate}

What survives must either break the involution symmetry in the weighting/domain, retain some
controlled non-invariant information, or prove a much stronger ensemble-specific structural
theorem.

\section{Obstruction IVd: Residual Symmetric Slices Still Stay At Chance}

The next natural reply is:
\begin{quote}
Fine.  We keep a little non-invariant information, so the exact sign-flip symmetry is broken.  We
do not need to resolve every orbit, only enough of them.
\end{quote}

This is the right next theorem to ask for.  The new Lean module
\texttt{ResidualSymmetryObstruction.lean} makes the burden precise: every residual slice on which
the retained information still fails to distinguish the involution partners remains exact-chance
under symmetry-respecting weights.

\subsection{Lean theorem}

\begin{theorem}[Any unresolved symmetric slice remains balanced]
Let \(\alpha\) be a finite sample space, let \(\tau : \alpha \to \alpha\) be an involution, let
\(p \subseteq \alpha\) be a \(\tau\)-stable slice, let \(u : \alpha \to U\) be the retained local
features, let \(y : \alpha \to \{0,1\}\) be the target bit, let \(w : \alpha \to M\) be a weight
into an additive commutative monoid, and let \(h : U \to \{0,1\}\) be any classifier.  Assume on
the slice \(p\) that
\[
  u(\tau x) = u(x), \qquad y(\tau x) = 1-y(x), \qquad w(\tau x) = w(x).
\]
Then the total weighted mass of correctly classified points in \(p\) equals the total weighted
mass of incorrectly classified points in \(p\).
\end{theorem}

\noindent
The key formal statements are:
\begin{itemize}[leftmargin=1.5em]
  \item \texttt{weightedCorrectMass}: total weighted correct mass for a feature-only classifier;
  \item \texttt{weightedIncorrectMass}: total weighted incorrect mass;
  \item \texttt{weightedCorrectMass\_eq\_weightedIncorrectMass}: the global weighted half-accuracy
  theorem under exact involution symmetry;
  \item \texttt{weightedCorrectMass\_eq\_weightedIncorrectMass\_on\_slice}: the residual-slice
  version obtained by restricting to a \(\tau\)-stable subtype.
\end{itemize}

\subsection{Why this hits the ``controlled non-invariant info'' reply}

This is the right quantitative obstruction.  A small amount of non-invariant information can only
help on the part of the domain where it actually breaks the orbit symmetry.  Every remaining
slice on which the feature map is still involution-invariant contributes only chance performance.

So the paper can no longer get by with a vague statement of the form:
\begin{quote}
We keep a little symmetry-breaking local information, and that should be enough to restore a
useful predictor.
\end{quote}

To make that route work, one now needs a real theorem that the unresolved symmetric slice carries
negligible weight, or else a theorem that the retained information breaks the involution on almost
all of the relevant mass.

\subsection{Consequence}

This does not yet show that \emph{every} controlled symmetry-breaking repair fails.  It does show
that such a repair only helps to the extent that it resolves most of the mass.  What remains must
therefore prove one of:
\begin{enumerate}[leftmargin=1.5em]
  \item the residual symmetric slice has negligible weight,
  \item the retained non-invariant information separates almost all involution pairs,
  \item or the actual evaluation domain is not well-modeled by a symmetry-respecting weighted
  slice at all.
\end{enumerate}

\section{Obstruction IVe: Advantage Is Bounded By Symmetry-Broken Mass}

The next creative rescue move is the quantitative version of the previous one:
\begin{quote}
Maybe we only break the symmetry on a relatively small region, but that is enough because the
classifier can exploit that region very efficiently and still gain a substantial global advantage.
\end{quote}

The new Lean module \texttt{AsymmetryBudgetObstruction.lean} blocks that move in the natural
finite weighted model.

\subsection{Lean theorem}

\begin{theorem}[Asymmetry budget bound]
Let \(\alpha\) be a finite sample space with weight \(w : \alpha \to \mathbb{N}\).  Let
\(p \subseteq \alpha\) be a \(\tau\)-stable slice on which the retained feature map is still
invariant under an involution \(\tau\), the target bit is flipped by \(\tau\), and the weights are
\(\tau\)-invariant.  Then for every classifier \(h\),
\[
  2 \cdot \mathrm{CorrectMass}(h)
  \;\le\;
  \mathrm{TotalMass} + \mathrm{OutsideMass}(p).
\]
Equivalently, the global excess over chance is bounded by the total mass outside the unresolved
symmetric slice.
\end{theorem}

\noindent
The key formal statements are:
\begin{itemize}[leftmargin=1.5em]
  \item \texttt{two\_mul\_correctSliceMass\_eq\_sliceMass}: on the unresolved symmetric slice, the
  correct mass is exactly half of the slice mass;
  \item \texttt{weightedCorrectMass\_eq\_correctSlice\_add\_correctOutside}: total correct mass
  splits into the unresolved slice contribution plus the symmetry-broken contribution;
  \item \texttt{correctOutsideMass\_le\_outsideMass}: outside the unresolved slice, even perfect
  classification cannot use more mass than is actually there;
  \item \texttt{two\_mul\_weightedCorrectMass\_le\_total\_plus\_outside}: the final asymmetry-budget
  inequality.
\end{itemize}

\subsection{Why this hits the last soft escape}

This is stronger than saying ``the residual slice is still at chance.''  It says the \emph{entire
global advantage} is quantitatively budgeted by the part of the domain where symmetry is actually
broken.

So one can no longer say:
\begin{quote}
We only need a little symmetry-breaking local information.  Perhaps it affects a small region, but
that small region could still be enough to drive the whole lower-bound argument.
\end{quote}

Not without a theorem that the symmetry-broken region itself carries substantial mass.  The Lean
bound shows that if the unresolved symmetric slice has mass close to the total, then the global
accuracy remains close to \(1/2\).

\subsection{Consequence}

At this point, the remaining route is no longer merely ``add some asymmetry.''  It must prove a
real mass-transfer statement:
\begin{enumerate}[leftmargin=1.5em]
  \item either almost all relevant mass leaves the involution-symmetric slice after the repair,
  \item or the weighted domain relevant to the proof is not symmetry-respecting in the sense used
  above,
  \item or the real predictor uses stronger ensemble-specific information than the current local
  symmetry analysis captures.
\end{enumerate}

\section{Obstruction IVf: Invariant Soft Scores Have Zero Signed Signal}

Another creative move is:
\begin{quote}
Perhaps the proof should not be phrased as conditional means on feature fibers at all.  Maybe one
can compute a richer soft score from the invariant local summary and then feed that score into a
more delicate downstream argument.
\end{quote}

This is a real strengthening of the repair attempt.  The new Lean module
\texttt{InvariantScoreObstruction.lean} blocks the core signal claim under the same involution
symmetry.

\subsection{Lean theorem}

\begin{theorem}[Invariant soft scores have zero signed correlation]
Let \(\alpha\) be a finite sample space, let \(\tau : \alpha \to \alpha\) be an involution, let
\(u : \alpha \to U\) be a retained feature map, let \(y : \alpha \to \{0,1\}\) be the target bit,
let \(w : \alpha \to \mathbb{N}\) be a \(\tau\)-invariant weight, and let
\(\mathrm{score} : U \to \mathbb{Z}\) be any soft score computed only from the retained features.
Assume
\[
  u(\tau x) = u(x), \qquad y(\tau x) = 1-y(x), \qquad w(\tau x) = w(x)
\]
for every \(x \in \alpha\).  Writing
\[
  \mathrm{sgnTarget}(y) =
  \begin{cases}
    1 & y=1, \\
    -1 & y=0,
  \end{cases}
\]
one has
\[
  \sum_{x \in \alpha} w(x)\,\mathrm{score}(u(x))\,\mathrm{sgnTarget}(y(x)) = 0.
\]
\end{theorem}

\noindent
The key formal statements are:
\begin{itemize}[leftmargin=1.5em]
  \item \texttt{targetSign}: the \(\{\pm 1\}\)-encoding of a Boolean target;
  \item \texttt{targetSign\_flip}: flipping the Boolean target negates that signed encoding;
  \item \texttt{weighted\_signedScore\_sum\_eq\_zero}: the weighted zero-correlation theorem for
  any invariant soft score.
\end{itemize}

\subsection{Why this hits the richer invariant-scoring repair}

This closes a subtler loophole than the earlier fiber-balance statements.  Those showed that the
conditional mean on each invariant feature value is exactly \(1/2\), and therefore that hard
classifiers or direct mean estimators on those features cannot beat chance.  The new theorem says
more: even before one thresholds or calibrates anything, \emph{the score itself} has zero signed
signal against the target.

So one can no longer say:
\begin{quote}
We are not really using a classifier on the invariant features.  We only need a good enough
numerical score from those features, and then a later nonlinear step can extract the useful
signal.
\end{quote}

Not under the same involution symmetry.  Any score depending only on the invariant retained
features is paired with its involution partner at equal weight and opposite signed target.  The
positive and negative contributions cancel exactly.

\subsection{Consequence}

Taken together with the previous sections, this leaves very little room for an invariant-feature
rescue.  The remaining route must now prove something outside the invariant scoring model:
\begin{enumerate}[leftmargin=1.5em]
  \item either the actual score retains genuinely non-invariant information,
  \item or the relevant domain/weighting is not symmetry-respecting,
  \item or the ensemble has a stronger special structure that invalidates the involution pairing
  hypotheses on the real predictor family.
\end{enumerate}

\section{Obstruction IVg: Invariant-Score Signal Comes Only From Weight Asymmetry}

The next and now very natural reply is:
\begin{quote}
Fine.  Perhaps the retained score inputs stay involution-invariant, but the real training or
evaluation weighting is not exactly symmetric.  Then an invariant soft score may still pick up a
useful signed signal.
\end{quote}

This is the last serious invariant-score escape hatch.  The new Lean module
\texttt{WeightAsymmetryObstruction.lean} makes that route exact: any invariant-score signal comes
entirely from the antisymmetric part of the weighting.

\subsection{Lean theorem}

\begin{theorem}[Only antisymmetric weight mass contributes to invariant-score signal]
Let \(\alpha\) be a finite sample space, let \(\tau : \alpha \to \alpha\) be an involution, let
\(u : \alpha \to U\) be a retained feature map, let \(y : \alpha \to \{0,1\}\) be the target bit,
let \(w : \alpha \to \mathbb{N}\) be an arbitrary weight, and let
\(\mathrm{score} : U \to \mathbb{Z}\) be any soft score on the retained features.  Assume
\[
  u(\tau x) = u(x), \qquad y(\tau x) = 1-y(x)
\]
for every \(x \in \alpha\).  Then
\[
  2 \sum_{x \in \alpha} w(x)\,\mathrm{score}(u(x))\,\mathrm{sgnTarget}(y(x))
  =
  \sum_{x \in \alpha} (w(x)-w(\tau x))\,\mathrm{score}(u(x))\,\mathrm{sgnTarget}(y(x)).
\]
\end{theorem}

\noindent
The key formal statements are:
\begin{itemize}[leftmargin=1.5em]
  \item \texttt{signedScoreContribution}: one point's contribution to the invariant signed-score
  sum;
  \item \texttt{antisymmetricWeightContribution}: the corresponding contribution of the
  weight-difference term \(w(x)-w(\tau x)\);
  \item \texttt{two\_mul\_signedScore\_sum\_eq\_antisymmetricWeight\_sum}: the exact decomposition
  identity.
\end{itemize}

\subsection{Why this narrows the weighting rescue}

This is stronger than the earlier exact-symmetry obstruction.  It says not just that symmetric
weights give zero signal, but that \emph{all} invariant-score signal is sourced only by the
antisymmetric part of the weighting.  The involution-invariant part cancels exactly.

So one can no longer say:
\begin{quote}
The true weighting is probably a little asymmetric, and that may be enough.
\end{quote}

Not without quantifying that asymmetry.  The theorem shows that the proof now needs a real bound of
the form:
\begin{quote}
the antisymmetric weight budget is large enough on the relevant domain to support the claimed
advantage.
\end{quote}

\subsection{Consequence}

This pushes the remaining route into a much narrower quantitative corner.  To save invariant local
scoring, one must now prove one of:
\begin{enumerate}[leftmargin=1.5em]
  \item the actual proof-relevant weighting has a substantial antisymmetric component,
  \item the score retains genuinely non-invariant information, or
  \item the real ensemble/predictor pair is not governed by the involution model above.
\end{enumerate}

\section{Obstruction IVh: Side Channels Help Only On Orbit-Resolving Mass}

At this point the natural pivot is:
\begin{quote}
Fine.  We stop asking invariant features alone to do the whole job.  We keep a controlled
involution-sensitive side channel \(v\), and classify on the pair \((u,v)\).
\end{quote}

This is the right next charitable target.  The new Lean module
\texttt{SideChannelResolutionObstruction.lean} packages the earlier symmetry budget results into the
exact theorem that this move needs to face.

\subsection{Lean theorem}

\begin{theorem}[Side-channel advantage is bounded by orbit-resolving mass]
Let \(\alpha\) be a finite sample space, let \(\tau : \alpha \to \alpha\) be an involution, let
\(u : \alpha \to U\) be a retained invariant feature map, let \(v : \alpha \to V\) be an arbitrary
side channel, let \(y : \alpha \to \{0,1\}\) be the target bit, let \(w : \alpha \to \mathbb{N}\)
be a \(\tau\)-invariant weight, and let \(h : U \times V \to \{0,1\}\) be any classifier.  Assume
\[
  u(\tau x) = u(x), \qquad y(\tau x) = 1-y(x), \qquad w(\tau x) = w(x)
\]
for every \(x \in \alpha\).  Define the \emph{resolved mass} to be the total weight of the points
for which the side channel separates the involution pair:
\[
  \mathrm{ResolvedMass}(v)
  =
  \sum_{x : v(\tau x) \neq v(x)} w(x).
\]
Then
\[
  2 \cdot \mathrm{CorrectMass}(h \circ (u,v))
  \;\le\;
  \mathrm{TotalMass} + \mathrm{ResolvedMass}(v).
\]
Equivalently, the entire global excess over chance is bounded by the weight of the domain where the
side channel actually resolves the orbit pairing.
\end{theorem}

\noindent
The key formal statements are:
\begin{itemize}[leftmargin=1.5em]
  \item \texttt{unresolvedBySideChannel}: the slice on which \(v(\tau x)=v(x)\);
  \item \texttt{resolvedMass}: the weight outside that unresolved slice;
  \item \texttt{two\_mul\_weightedCorrectMass\_pair\_le\_total\_plus\_resolvedMass}: the final
  resolution-budget inequality for classifiers on \((u,v)\).
\end{itemize}

\subsection{Why this hits the controlled-side-channel rescue}

This is stronger than the earlier residual-symmetry prose, because it says exactly what the side
channel must buy.  Keeping some involution-sensitive information does \emph{not} by itself create a
large advantage.  The benefit is quantitatively capped by the mass where that side channel really
separates involution partners.

So one can no longer say:
\begin{quote}
We only need a small side channel.  Once it is present, a clever classifier on \((u,v)\) may still
recover a substantial domination advantage.
\end{quote}

Not without a theorem that the side channel resolves involution pairs on a substantial fraction of
the relevant mass.  The unresolved slice remains chance, and the total excess over chance is budgeted
by what lies outside it.

\subsection{Consequence}

This leaves the controlled-side-channel route in a much sharper form.  To save it, one must now
prove one of:
\begin{enumerate}[leftmargin=1.5em]
  \item the retained side channel separates involution pairs on most of the proof-relevant mass,
  \item the weighting itself has a large antisymmetric component that can be exploited,
  \item or the real predictor is not well-modeled by this involution-paired finite setup.
\end{enumerate}

\section{Obstruction IVi: Any Domination Gap Forces Matching Resolution}

The previous section bounded success by resolved mass.  The converse reading is now just as
important:
\begin{quote}
if the proof claims a real domination gap over chance from a retained side channel, that claim
already commits it to a quantitatively comparable orbit-resolution theorem.
\end{quote}

The new Lean module \texttt{ResolutionDemandObstruction.lean} packages that commitment as a direct
corollary.

\subsection{Lean theorem}

\begin{theorem}[Doubled advantage is bounded by resolved mass]
For a weighted classifier \(h\), define the \emph{doubled advantage}
\[
  \mathrm{DoubledAdv}(h)
  =
  2 \cdot \mathrm{CorrectMass}(h) - \mathrm{TotalMass}.
\]
Under the same involution hypotheses as in Obstruction IVh, one has
\[
  \mathrm{DoubledAdv}(h \circ (u,v))
  \le
  \mathrm{ResolvedMass}(v).
\]
\end{theorem}

\noindent
The key formal statements are:
\begin{itemize}[leftmargin=1.5em]
  \item \texttt{doubledAdvantage}: twice the weighted success gap over chance;
  \item \texttt{doubledAdvantage\_pair\_le\_resolvedMass}: the final converse bound.
\end{itemize}

\subsection{Why this matters}

This is the strongest current version of the side-channel blocker.  It no longer says only that
resolved mass is \emph{sufficient} for advantage.  It says a claimed advantage is itself evidence
that substantial resolution must already be present.

So one can no longer say:
\begin{quote}
Perhaps the predictor gets a modest but important edge from a very small side channel, and only
later amplifies that edge through clever structure.
\end{quote}

Not unless that modest edge is matched by actual orbit-resolving mass.  The theorem makes the
burden explicit: any nontrivial domination statement already entails a nontrivial local resolution
statement.

\subsection{Consequence}

This is close to the form needed to challenge the route on its own terms.  To save the current
style of proof, one now needs a theorem of the form:
\begin{quote}
the switched predictor's retained side channel resolves involution pairs on enough mass to support
the claimed lower-bound-driving advantage.
\end{quote}

Without that, the claimed domination gap has no quantitative room to exist.

\section{Obstruction V: Short Programs Still Cover the Full Log-Width Local Rule Class}

Another natural reply is:
\begin{quote}
The switched predictor comes from a short global program, so even if arbitrary local rules are
huge, only a much smaller subclass can actually arise from polynomial-time decoders.
\end{quote}

This is not true without an additional restriction beyond mere global shortness.

\subsection{Lean theorem}

The new module \texttt{ShortProgramObstruction.lean} formalizes the basic truth-table fact.

\begin{theorem}[Truth tables realize all local rules]
For every \(n\), there is a surjective decoder
\[
  \{0,1\}^{2^n} \twoheadrightarrow \{\,f : \{0,1\}^n \to \{0,1\}\,\}.
\]
Equivalently, every Boolean rule on \(n\) visible bits can be hard-wired by a truth table of
length \(2^n\).
\end{theorem}

\noindent
The key formal statements are:
\begin{itemize}[leftmargin=1.5em]
  \item \texttt{truthTableDecode\_surjective}: every local rule is realized by a
  \(2^n\)-bit truth table;
  \item \texttt{truthTableFamily\_realizes\_all}: the realized class of the truth-table family is
  the full local-rule space;
  \item \texttt{card\_realized\_truthTableFamily}: this realized class has exactly cardinality
  \(2^{2^n}\).
\end{itemize}

\subsection{Binary-log consequence}

At binary-log width the truth-table size is already linear in the ambient parameter:
\[
  2^{\lfloor \log_2 m \rfloor + 1} \le 2m \qquad (m \ne 0).
\]
This is formalized as
\texttt{truthTableBits\_binaryLogWidth\_le\_twoMul}.

So if the switched interface really has only \(n = \lfloor \log_2 m \rfloor + 1\) visible bits,
then a merely linear-size code budget already suffices to realize the \emph{entire} local-rule
class on that interface.

\subsection{Consequence}

This directly blocks the short-program rescue in its naive form:

\begin{quote}
``The global decoder is polynomial-size'' does not imply
``the induced local rule lies in a tiny subclass.''
\end{quote}

On the contrary, polynomial-size descriptions are already large enough to hard-wire every Boolean
rule on \(O(\log m)\) visible bits.  So if one wants a true small subclass theorem, it must use
something much stronger than global shortness alone: symmetry, algebraic normal form, a tree
interpreter with bounded message alphabet, or some other genuine structural restriction.

\section{Obstruction VI: Exact Metagraph Quotients Have Interface-State Explosion}

The strongest remaining charitable rescue is roughly this:
\begin{quote}
The switched neighborhood first collapses to a canonical metagraph or message-passing quotient,
and the witness bit is computed by a fixed interpreter on that quotient.
\end{quote}

This is a serious idea.  It is also constrained by a generic exactness barrier.

\subsection{Lean theorem}

The new module \texttt{MetagraphObstruction.lean} models a quotient state by its exact interface
behavior on \(k\) visible boundary bits.  The semantic object is then a boundary transducer
\[
  (\{0,1\}^k \to \{0,1\}^k).
\]

\begin{theorem}[Boundary transducers are exponentially numerous]
There are exactly
\[
  2^{k 2^k}
\]
boundary transducers on \(k\) visible bits.
\end{theorem}

\noindent
The key formal statements are:
\begin{itemize}[leftmargin=1.5em]
  \item \texttt{card\_boundaryMap}: exact \(k\)-bit interface semantics have cardinality
  \(2^{k2^k}\);
  \item \texttt{no\_surjective\_bitCode\_to\_boundaryMap\_of\_lt}: no \(s\)-bit code can cover the
  full interface-semantic space when \(s < k2^k\);
  \item \texttt{no\_exact\_metagraphQuotient\_below\_inputSize\_binaryLogWidth}: at
  \(k = \lfloor \log_2 m \rfloor + 1\), there is no exact quotient with fewer than \(m\) code bits.
\end{itemize}

\subsection{Why this matters}

If a metagraph/message quotient is meant to be \emph{exact}, then it must preserve enough state to
reconstruct the gadget's interface behavior under gluing.  The theorem above says that, in the
absence of an additional structural restriction, the exact interface-semantic space is already far
too large for a tiny finite-state algebra.

This does not yet show that the current ensemble realizes the full transducer class.  What it
does show is that a small exact quotient theorem cannot follow from tree-likeness or locality
alone.  To survive, the metagraph route needs a new theorem that the actual switched/VV gadgets
lie in a drastically smaller semantic subclass.

\subsection{Consequence}

So the metagraph rescue is now narrowed to something like:
\begin{quote}
not ``there is an exact small quotient of all local interfaces,''
but ``the specific random masked+isolated interfaces collapse to a much smaller algebra for
ensemble-specific reasons.''
\end{quote}

That is a legitimate possible theorem.  It is also substantially stronger than anything currently
written in the paper, and it is the next point where a decisive break may occur.

\section{Obstruction VII: Tree-Likeness Alone Does Not Yield a Small Message Class}

Another natural repair is:
\begin{quote}
On locally tree-like neighborhoods, the witness bit is computed by a small message-passing or
dynamic-program interpreter on the tree.
\end{quote}

This too needs a stronger premise than mere tree structure.

\subsection{Lean theorem}

The new module \texttt{TreeMessageObstruction.lean} defines a fixed perfect binary tree
architecture of depth \(n\), where internal nodes merely branch on the next input bit and only the
leaves carry data.

\begin{theorem}[Perfect binary trees realize all local rules]
For every Boolean rule
\[
  f : \{0,1\}^n \to \{0,1\},
\]
there is a depth-\(n\) perfect binary decision tree whose evaluation equals \(f\).
\end{theorem}

\noindent
The key formal statements are:
\begin{itemize}[leftmargin=1.5em]
  \item \texttt{eval\_encodeDecisionTree}: every local rule is exactly realized by a leaf-labeled
  depth-\(n\) perfect tree;
  \item \texttt{evalDecisionTree\_surjective}: the semantics map from such trees onto
  \texttt{LocalRule n} is surjective;
  \item \texttt{no\_surjective\_bitCode\_to\_decisionTreeSemantics\_of\_lt}: therefore no
  \(s\)-bit exact code can cover all depth-\(n\) tree semantics when \(s < 2^n\).
\end{itemize}

\subsection{Why this matters}

This is a direct answer to the naive tree-message rescue:

\begin{quote}
tree-likeness by itself does \emph{not} imply semantic simplicity.
\end{quote}

Even with an extremely rigid tree architecture, the semantic class is already the full Boolean
function space on the visible inputs.  So if one wants a genuine tree-message theorem, it must use
special properties of the actual SAT+VV messages, not just the fact that the underlying graph is a
tree.

\subsection{Consequence}

The surviving charitable tree route is now much narrower:
\begin{quote}
not ``all tree-like local neighborhoods admit a tiny interpreter,''
but ``the specific masked+isolated tree messages belong to a sharply smaller algebra.''
\end{quote}

That is again a real possible theorem.  It is also much stronger than the current text, and it is
precisely the kind of statement that now has to be made explicit to keep this proof approach alive.

\section{What These Results Do Not Yet Prove}

The current Lean modules do \emph{not} prove that the full argument fails.  In particular,
they do not yet address:
\begin{enumerate}[leftmargin=1.5em]
  \item whether the masking and isolation ensemble enforces extra algebraic structure on local rules;
  \item whether the switched predictors might collapse to a subclass analogous to low-depth message
  passing or tree computations;
  \item whether the neutrality and sparsification lemmas are themselves correct in the precise form
  used in the paper;
  \item whether a different, stronger switching theorem could bypass the present obstruction.
\end{enumerate}

So the current status is best described as:
\begin{quote}
The proof attempt is not yet broken in full, but one of its central complexity-compression
transitions has not been justified, and the default combinatorics point in the opposite direction.
\end{quote}

\section{Next Attack Surface}

The next natural formal target is the paper's implicit compression claim.  The
first and third interface steps are now isolated in Lean:
\texttt{VisiblePostSwitchSurface.lean} defines the exact manuscript-visible
surface \(u=(z,a,b)\) together with its invariant and fork projections, and
\texttt{ExactSwitchedFamily.lean} packages the exact bit-budget theorem needed
to reuse the same-route ERM bound.  The remaining mathematical burden is now
sharper:

\begin{quote}
After switching, the witness-bit predictor is not merely local; it is simple enough to lie in a
small explicit class.
\end{quote}

To make that rigorous, one would want a theorem of the shape:
\begin{enumerate}[leftmargin=1.5em]
  \item define the exact visible post-switch data carried by the masked Unique-SAT ensemble;
  \item define the induced witness-bit function on that data;
  \item prove that this function belongs to a sharply bounded class
  (for example, bounded-depth circuits of size \(\mathrm{polylog}(m)\), or a code family of
  description length \(\mathrm{polylog}(m)\));
  \item verify that this bound is strong enough to feed the later union-bound and compression steps.
\end{enumerate}

The new module \texttt{ExactVisibleCompressionObstruction.lean} records the
fallback counting barrier on that exact object: if the switched family on
\(u=(z,a,b)\) is still rich enough to realize \emph{all} Boolean rules on the
exact visible surface, then no \(s\)-bit code can cover it when
\(s < |\text{surface}|\).  So any same-route rescue must now prove a genuine
strict subclass theorem for the \emph{actual} switched predictors.

If such a theorem cannot be proved, then the proof attempt should be regarded as stalled at a real
mathematical gap, not a missing editorial detail.

\section{Lean Excerpts}

\begin{lstlisting}[caption={Core asymptotic obstruction}]
theorem polylog_isLittleO_logRadiusNeighborhood (k : Nat) {d c : Real}
    (hd : 1 < d) (hc : 0 < c) :
    polylog k =o[atTop] logRadiusNeighborhood d c

theorem logRadiusNeighborhood_not_isBigO_polylog (k : Nat) {d c : Real}
    (hd : 1 < d) (hc : 0 < c) :
    ¬ logRadiusNeighborhood d c =O[atTop] polylog k
\end{lstlisting}

\begin{lstlisting}[caption={Core counting obstruction}]
theorem no_injective_bitCode_of_lt {n s : Nat} (hs : s < 2 ^ n) :
    ¬ ∃ f : LocalRule n → BitCode s, Function.Injective f

theorem no_uniform_code_below_expInput_for_binaryLogNeighborhood
    {d m s : Nat} (hd : 2 ≤ d) (hs : s < 2 ^ m) :
    ¬ ∃ f : LocalRule (d ^ (Nat.log 2 m + 1)) → BitCode s,
      Function.Injective f
\end{lstlisting}

\section{Conclusion}

The present crux-first formal work yields a clean conditional diagnosis.

\begin{enumerate}[leftmargin=1.5em]
  \item Logarithmic-radius locality is \emph{not} the same thing as polylogarithmic complexity.
  \item The full class of local Boolean rules on that visible data is far too large to fit into a
  polylogarithmic coding budget.
  \item Therefore the program needs a separate theorem proving exceptional compressibility
  of the particular switched predictors it generates.
\end{enumerate}

That is currently the sharpest formally checked crux we have in Lean.

\begin{thebibliography}{9}
\bibitem{manuscript2025}
Anonymous manuscript under study.
\newblock \emph{P != NP: A Non-Relativizing Proof via Quantale Weakness and Geometric Complexity}.
\newblock October 2025.

\bibitem{mettapedia-pnp}
Mettapedia formal notes.
\newblock \texttt{Mettapedia/Computability/PNP/LocalityObstruction.lean} and
\texttt{Mettapedia/Computability/PNP/CompressionObstruction.lean}.
\end{thebibliography}

\end{document}
